\section{System Architecture}
\label{sec:architecture}

\subsection{Overview}

OpenSkynet is organized into four conceptual layers, as illustrated in Figure~\ref{fig:architecture}: (1) the \textbf{orchestration layer}, containing the Manager Agent, Agent Loop, and subagent system; (2) the \textbf{skill layer}, responsible for learning, executing, healing, and auditing browser automation skills; (3) the \textbf{memory layer}, providing tiered persistent storage with hybrid retrieval; and (4) the \textbf{browser layer}, abstracting over multiple execution backends.

\begin{figure*}[t]
    \centering
    \setlength{\fboxsep}{0pt}%
    \begin{minipage}{\textheight}
    \begin{tikzpicture}[
        x=2.2cm, y=1.8cm,
        every node/.style={font=\sffamily, rounded corners=4pt},
        sm/.style={
            minimum height=0.9cm, minimum width=2.0cm,
            text=white, font=\sffamily\tiny, align=center,
            fill, draw=none,
        },
        lg/.style={
            minimum height=2.0cm, minimum width=2.4cm,
            text=white, font=\sffamily\scriptsize, align=center,
            fill, draw=none,
        },
        arr/.style={-{Stealth[length=3pt, width=2pt]}, color=cArrow, very thick},
        lbl/.style={font=\sffamily\tiny, text=cArrow, fill=white, inner sep=1pt},
        be/.style={
            draw=cBrowser!40, rounded corners=3pt, thick,
            fill=cBrowser!10,
            minimum height=0.6cm, minimum width=2.6cm,
            text=cBorder, font=\sffamily\scriptsize, align=center,
        },
    ]
        % TOP ROW (y=4): User -- Manager -- LLM
        \node[sm, fill=cInput]    at (0,4)   (u) {\faUser\\\textbf{User Input}\\[-1pt]CLI/TUI/API};
        \node[sm, fill=cManager]  at (3,4)   (m) {\faSitemap\\\textbf{Manager Agent}\\[-1pt]Strategy \,\textbullet\, Decompose};
        \node[sm, fill=cLLM]      at (6,4)   (l) {\faBrain\\\textbf{LLM Providers}\\[-1pt]OpenAI \,\textbullet\, Ollama};

        % CONNECT: top row
        \draw[arr] (u.east)  -- (m.west);
        \draw[arr] (m.east) -- node[lbl, above] {planning} (l.west);
        \draw[arr] (m.south) -- ++(0,-0.4) -- (2.0, 2.6) node[lbl, right] {plan};

        % MIDDLE ROW (y=2): Skills | Loop | Subagents | Memory
        \node[lg, fill=cSkills]    at (0,2)   (sk) {%
            \faCogs\\\textbf{Skill Engine}\\[3pt]
            \faGraduationCap~Learn\quad\faMedkit~Heal\\[0pt]
            \faSearch~Search\quad\faClipboardCheck~Audit};

        \node[lg, fill=cAgent]      at (3,2)   (loop) {%
            \faSync\\\textbf{Agent Loop}\\[3pt]
            \faClipboardList~Plan\quad\faPlay~Execute\\[0pt]
            \faEye~Observe\quad\faLightbulb~Reflect};

        \node[lg, fill=cSubagent]   at (6,2)   (sub) {%
            \faUsersCog\\\textbf{Subagents}\\[3pt]
            \faGlobe~browser\quad\faCode~code\\[0pt]
            \faBug~debug\quad\faSearchPlus~explore\\[0pt]
            \faUserShield~redteam};

        \node[lg, fill=cMemory]    at (9,2)   (mem) {%
            \faDatabase\\\textbf{Tiered Memory}\\[3pt]
            {Working $\to$ Session $\to$ Long-Term}\\[0pt]
            \faVectorSquare~Vector\quad\faRecycle~Consolidate};

        % CONNECT middle row
        \draw[arr] (loop.west) -- node[lbl, above] {learn/heal} (sk.east);
        \draw[arr] (loop.east) -- node[lbl, above] {delegate} (sub.west);
        \draw[arr] (loop.north east) -- ++(0.3,0) |- node[lbl, near start] {store} (mem.west);
        \draw[arr] (mem.north west) -- ++(-0.3,0) |- node[lbl, near start] {context} (loop.north east);

        % BOTTOM ROW (y=0): Browser layer
        \node[draw=cBrowser!30, rounded corners=5pt, thick,
              fill=cBrowser!6,
              minimum height=1.0cm, minimum width=12cm,
              text=cBrowser, font=\sffamily\scriptsize\bfseries] at (4.5,0) (bl)
            {\faWindowMaximize~~Browser Automation Layer};

        \node[be] at (1.0,-0.35) (b1) {\faDesktop~~Playwright};
        \node[be] at (4.5,-0.35) (b2) {\faNodeJs~~Agent-Browser};
        \node[be] at (8.0,-0.35) (b3) {\faCog~~OpenBrowser (Rust)};

        \draw[arr] (sk.south) -- ++(0,-0.3) -| (b1.north);
        \draw[arr] (sub.south) -- ++(0,-0.3) -| (b2.north);
        \draw[arr] (mem.south) -- ++(0,-0.3) -| (b3.north);
        \draw[arr] (loop.south) -- ++(0,-1.0);

        % LEGEND
        \node[anchor=north west, font=\sffamily\tiny, text=cBorder, inner sep=0pt]
            at (-1, -0.9) {%
            \color{cManager}\rule{4pt}{4pt}\;Orchestration\quad
            \color{cSkills}\rule{4pt}{4pt}\;Skills\quad
            \color{cMemory}\rule{4pt}{4pt}\;Memory\quad
            \color{cBrowser}\rule{4pt}{4pt}\;Browser\quad
            \color{cSubagent}\rule{4pt}{4pt}\;Subagents};
    \end{tikzpicture}
    \end{minipage}
    \caption{OpenSkynet system architecture.}
    \label{fig:architecture}
\end{figure*}

\subsection{Manager Agent and Strategy Selection}

The Manager Agent~\cite{browseruse} is the entry point for all user requests. Upon receiving a task, it selects one of five strategies based on the task type and conversation context:

\begin{itemize}[leftmargin=*,itemsep=1pt]
    \item \textbf{Conversational.} For informational queries that require no browser interaction, the manager responds directly.
    \item \textbf{Direct.} A straightforward browser automation task is passed directly to the Browser Agent.
    \item \textbf{Use-Skill.} If a learned skill matches the task description (via vector or keyword similarity), the existing skill is retrieved and executed.
    \item \textbf{Delegate.} Coding tasks are routed to the code subagent; complex tasks may be dispatched to specialized subagents.
    \item \textbf{Decompose.} Multi-phase tasks are decomposed into independent subtasks using beam-width planning (up to 5 subtasks across 2 beam candidates), scored by task coverage and parallelizability.
\end{itemize}

Strategy selection uses a multi-tier fast path: explicit URL navigation bypasses the LLM entirely; keyword-matching identifies coding tasks; an LLM classifier resolves ambiguous cases; and the full LLM planner is invoked only when necessary. This minimizes token consumption for simple tasks while preserving the model's reasoning capability for complex ones.

\subsection{Agent Loop}

The Agent Loop implements a Plan--Execute--Observe--Reflect cycle inspired by the ReAct paradigm~\cite{react}. At each iteration:

\begin{enumerate}[leftmargin=*,itemsep=1pt]
    \item \textbf{Plan:} The manager produces a plan with strategy tags and optional subtask decomposition.
    \item \textbf{Execute:} The current step is dispatched according to its strategy (direct, skill-based, or delegated to a subagent).
    \item \textbf{Observe:} The result is classified as success or failure; screenshots and DOM snapshots are captured for downstream analysis.
    \item \textbf{Reflect:} An LLM-based reflection agent evaluates whether the task is complete, with confidence scoring. If confidence is low or the step failed, a retry with exponential backoff ($2^n + \epsilon$, capped at 10s) is attempted. After exhausting retries, the system falls back through strategies: Use-Skill $\rightarrow$ Direct, Delegate $\rightarrow$ Direct, Decompose $\rightarrow$ Delegate.
\end{enumerate}

A Budget tracker monitors replans (max 3), total retries (max 8), wall-clock time (max 600s), LLM tokens (max 200K), and browser actions (max 100). When any resource is exhausted, the agent terminates gracefully.

\subsection{Prompt Optimization}

To reduce token consumption, OpenSkynet provides three system prompt tiers:

\begin{itemize}[leftmargin=*,itemsep=1pt]
    \item \textbf{Full mode} (default): Comprehensive system prompt including agent identity, memory context, project context, and task-adaptive add-ons.
    \item \textbf{Flash mode}: Compact system prompt retaining verification rules and error recovery instructions but omitting persona and context.
    \item \textbf{Turbo mode}: Ultra-minimal prompt (3 lines) for simple, well-defined tasks with zero context overhead.
\end{itemize}

Task-adaptive add-ons are dynamically appended based on keyword matching: extraction tasks receive structured output instructions; form-filling tasks receive field-by-field guidance; search tasks receive query construction strategies; and navigation tasks receive multi-page state management rules.

\subsection{Browser Abstraction Layer}

OpenSkynet supports three browser backends through a unified interface:

\begin{enumerate}[leftmargin=*,itemsep=1pt]
    \item \textbf{Playwright / Browser-Use:} The primary backend, providing full Playwright control with stealth mode (CloakBrowser), cookie/localStorage persistence, and 18 discrete browser tools exposed to the LLM (navigate, click, type, scroll, screenshot, extract, etc.).
    \item \textbf{Agent-Browser Sidecar:} A Node.js process that runs a self-contained AI agent with its own browser control. Communication occurs via JSON-RPC 2.0 over stdin/stdout. Useful when the main agent delegates browsing to an isolated subagent.
    \item \textbf{OpenBrowser REST:} A Rust-based semantic browser that exposes a REST API with semantic tree introspection. The controller adapter translates REST responses into the same \texttt{PageSnapshot} and \texttt{ElementInfo} data structures used by the Playwright backend.
\end{enumerate}

All backends share a common \texttt{BrowserSession} interface (start/stop/screenshot/save-state/load-state) and a common \texttt{BrowserController} interface (18 browser actions). A module-level singleton controller enables transparent backend swapping without modifying tool definitions.
